\documentclass[conference]{IEEEtran}
\IEEEoverridecommandlockouts

\usepackage{cite}
\usepackage{amsmath,amssymb,amsfonts}
\usepackage{algorithmic}
\usepackage{graphicx}
\usepackage{textcomp}
\usepackage{xcolor}
\usepackage{booktabs}
\usepackage{array}
\usepackage{tabularx}
\usepackage{hyperref}
\usepackage{listings}

\lstset{
	basicstyle=\ttfamily\footnotesize,
	breaklines=true,
	frame=single,
	backgroundcolor=\color{gray!10},
	keywordstyle=\color{blue},
	stringstyle=\color{red!70!black},
	commentstyle=\color{green!50!black},
	language=,
}

\def\BibTeX{{\rm B\kern-.05em{\sc i\kern-.025em b}\kern-.08em
		T\kern-.1667em\lower.7ex\hbox{E}\kern-.125emX}}

\begin{document}
	
	% ─────────────────────────────────────────────────────────────────────────────
	%  TITLE & AUTHORS
	% ─────────────────────────────────────────────────────────────────────────────
	\title{IoT-Based Multi-Node Firefighter Monitoring System\\
		with Edge AI and Real-Time Dashboard}
	
	\author{
		\IEEEauthorblockN{Rahul Nerella}
		\IEEEauthorblockA{
			\textit{Reg No: 24BEC0071}\\
			\textit{Dept. of Electronics and Communication Engineering}\\
			\textit{Vellore Institute of Technology (VIT)}\\
			Academic Year: 2025--2026
		}
		\and
		\IEEEauthorblockN{Sachith Velavan}
		\IEEEauthorblockA{
			\textit{Reg No: 24BEC0736}\\
			\textit{Dept. of Electronics and Communication Engineering}\\
			\textit{Vellore Institute of Technology (VIT)}\\
			Academic Year: 2025--2026
		}
	}
	
	\maketitle
	
	% ─────────────────────────────────────────────────────────────────────────────
	%  ABSTRACT
	% ─────────────────────────────────────────────────────────────────────────────
	\begin{abstract}
		Firefighting environments are highly hazardous, requiring real-time monitoring
		of personnel to ensure safety and operational efficiency. This paper presents a
		complete IoT-based multi-node firefighter monitoring system integrating embedded
		devices, edge-based TinyML inference, LoRa communication, and a real-time web
		dashboard. Each firefighter is equipped with an ESP32-based wearable device that
		collects environmental and motion data, performs local inference such as fall
		detection, air quality classification, and activity recognition, and transmits
		processed data via LoRa. A receiver node forwards the data to a Node.js server,
		which handles storage, parsing, and real-time visualization through a web-based
		dashboard. The system supports dynamic multi-node monitoring, emergency alerts,
		and efficient data logging, providing a scalable and reliable solution for
		firefighter safety.
	\end{abstract}
	
	\begin{IEEEkeywords}
		IoT, Firefighter Safety, TinyML, LoRa Communication, Real-Time Monitoring,
		Edge Computing, Embedded Systems.
	\end{IEEEkeywords}
	
	% ─────────────────────────────────────────────────────────────────────────────
	%  I. INTRODUCTION
	% ─────────────────────────────────────────────────────────────────────────────
	\section{Introduction}
	
	Firefighters operate in extreme conditions involving high temperatures, toxic
	gases, and structural instability. Continuous monitoring of their health and
	surroundings is critical for ensuring safety. Traditional systems lack real-time
	analytics and scalable monitoring capabilities. Existing solutions focus on
	either health monitoring or location tracking---not an integrated approach.
	Additionally, there is a lack of low-power long-range communication (LoRa)
	solutions specifically adapted for indoor fire scenarios, and limited real-time
	alert mechanisms for emergency conditions.
	
	This work proposes an IoT-based system that integrates edge intelligence,
	wireless communication, and real-time visualization to enhance firefighter
	safety and decision-making. The system is relevant to embedded systems, IoT,
	and communication (ECE domain) and offers novelty through the integration of
	GPS, health sensors, and LoRa communication using the ESP32 microcontroller.
	
	% ─────────────────────────────────────────────────────────────────────────────
	%  II. LITERATURE REVIEW
	% ─────────────────────────────────────────────────────────────────────────────
	\section{Literature Review}
	
	Table~\ref{tab:existing_work} summarizes the current state of work in the
	domain of firefighter monitoring systems.
	
	\begin{table}[h]
		\caption{Summary of Existing Work}
		\label{tab:existing_work}
		\centering
		\begin{tabular}{|l|p{5cm}|}
			\hline
			\textbf{Area} & \textbf{Existing Work} \\
			\hline
			Health Monitoring  & Wearable sensors for heart rate \& temperature \\
			\hline
			Location Tracking  & GPS-based outdoor tracking systems \\
			\hline
			Communication      & Wi-Fi / Zigbee for short-range transmission \\
			\hline
		\end{tabular}
	\end{table}
	
	Despite progress in individual areas, significant gaps remain:
	\begin{itemize}
		\item No unified system combining health monitoring, location tracking,
		and communication.
		\item GPS fails indoors; limited alternatives exist.
		\item Lack of LoRa-based firefighter monitoring systems.
		\item Wi-Fi fails in disaster zones due to short range and high power usage.
		\item No real-time emergency trigger mechanisms in prior work.
	\end{itemize}
	
	Methods used by prior studies include cloud-based IoT systems, Wi-Fi/Bluetooth
	communication, and basic environmental sensors. Their key limitations were short
	communication range, high power usage unsuitable for long missions, and the
	absence of real-time alert triggers.
	
	% ─────────────────────────────────────────────────────────────────────────────
	%  III. SYSTEM ARCHITECTURE
	% ─────────────────────────────────────────────────────────────────────────────
	\section{System Architecture}
	
	The proposed system is structured into three main layers: (1)~the Firefighter
	Node, (2)~the Communication Layer, and (3)~the Backend and Data Visualization
	Layer. This layered architecture ensures modularity, scalability, and efficient
	real-time operation.
	
	\subsection{The Firefighter Node}
	
	Each firefighter is equipped with an ESP32-based wearable device integrated
	with environmental sensors and GPS. The ESP32 microcontroller forms the backbone
	of the edge layer. It is a low-cost, low-power SoC with integrated Wi-Fi,
	Bluetooth, and multiple peripheral interfaces, making it suitable for embedded
	IoT systems~\cite{yadav2026}. The device performs local TinyML inference and
	transmits processed data via LoRa.
	
	\begin{table}[h]
		\caption{Experimental Component Setup}
		\label{tab:components}
		\centering
		\renewcommand{\arraystretch}{1.2}
		\begin{tabular}{|p{2.2cm}|p{2.4cm}|p{2.8cm}|}
			\hline
			\textbf{Component} & \textbf{Purpose} & \textbf{Reason for Selection} \\
			\hline
			ESP32 & Main controller & Low power, built-in WiFi/Bluetooth \\
			\hline
			LoRa Module (SX1278) & Long-range communication & Works in low-signal environments \\
			\hline
			GPS Module (NEO6M) & Location tracking & Accurate outdoor positioning \\
			\hline
			DHT22 & Temp/humidity monitoring & $\pm$0.5$^\circ$C accuracy, easy integration \\
			\hline
			MQ-2, MQ-9, MQ-135 & Toxic gas detection & Smoke, CO, and air pollutant detection \\
			\hline
			MAX30102 & Heart rate monitoring & Health tracking \\
			\hline
			MPU6050 IMU & Motion tracking & Fall detection and activity recognition \\
			\hline
			Flame Sensor & Fire detection & IR-based, complements gas/temp sensors \\
			\hline
		\end{tabular}
	\end{table}
	
	\noindent\textbf{Sender ESP32:}
	\begin{itemize}
		\item Reads all sensor data
		\item Performs TinyML inference (fall detection, activity recognition,
		air quality classification)
		\item Packages results into JSON format
		\item Transmits via LoRa
	\end{itemize}
	
	\noindent\textbf{Receiver ESP32:}
	\begin{itemize}
		\item Receives LoRa packets
		\item Forwards data to the backend server via serial communication
	\end{itemize}
	
	The MQ sensor series detect gases based on changes in resistance due to gas
	concentration. MQ-2 detects smoke and combustible gases; MQ-9 detects carbon
	monoxide (CO); MQ-135 detects air pollutants and toxic gases. Research shows
	that combining multiple MQ sensors improves environmental classification
	accuracy, especially when used with machine learning techniques~\cite{wang2024}.
	
	The MPU6050 provides 3-axis accelerometer and 3-axis gyroscope data. Research
	has demonstrated that such data can classify human activities such as walking,
	standing, and sitting with accuracy above 90\%~\cite{doddabasappla2021}.
	
	\subsection{Communication Layer (LoRa)}
	
	LoRa (Long Range) is a wireless communication technology designed for
	low-power wide-area networks (LPWANs), enabling communication over long
	distances with minimal energy consumption. LoRa is widely adopted in IoT
	systems due to its cost-effectiveness, scalability, and ability to operate in
	unlicensed frequency bands~\cite{haif2024}.
	
	LoRa allows communication exceeding 1~km in practical deployments while
	maintaining low energy consumption, making it ideal for emergency
	scenarios~\cite{obrien2025}. Advanced relaying techniques can increase coverage
	by up to 40\% and improve throughput and reliability~\cite{haif2024}.
	
	The system follows a star topology:
	\begin{itemize}
		\item Firefighter wearable devices act as LoRa end nodes
		\item Receiver ESP32 acts as the gateway
		\item Node.js server processes the forwarded data
	\end{itemize}
	
	Challenges such as packet collisions, limited data rate, and signal attenuation
	are mitigated using adaptive spreading factor allocation, frequency channel
	selection, and relay-based communication.
	
	\subsection{Backend and Visualization Layer}
	
	A Node.js server processes incoming data, stores it in multiple formats (JSON,
	CSV, SQLite), and transmits updates to a real-time dashboard using Socket.io.
	Node.js uses a single-threaded, event-driven architecture with non-blocking I/O,
	allowing it to handle thousands of concurrent connections
	efficiently~\cite{huang2020}.
	
	The backend follows a modular MVC architecture:
	\begin{itemize}
		\item \textbf{Model:} Stores firefighter data (JSON, CSV, SQLite)
		\item \textbf{Controller:} Processes incoming LoRa data
		\item \textbf{Route:} Handles API endpoints (GET, POST, PUT, DELETE)
	\end{itemize}
	
	WebSocket-based communication via Socket.io ensures real-time frontend updates
	without polling~\cite{chaplia2023}. Node.js supports microservices and
	distributed architectures, enabling the backend to scale to hundreds of
	firefighter nodes~\cite{chaplia2023}.
	
	% ─────────────────────────────────────────────────────────────────────────────
	%  IV. ALGORITHMS & EDGE INTELLIGENCE
	% ─────────────────────────────────────────────────────────────────────────────
	\section{Algorithms and Edge Intelligence}
	
	The system employs the following algorithms and inference models:
	
	\begin{itemize}
		\item \textbf{Threshold-based alert system} for temperature, gas, and heart rate.
		\item \textbf{GPS coordinate extraction} via NMEA parsing using TinyGPS++.
		\item \textbf{LoRa packet transmission protocol} with JSON-formatted payloads.
		\item \textbf{Sensor data filtering} using basic averaging across trial runs.
		\item \textbf{Emergency trigger logic} via hardware button interrupt.
	\end{itemize}
	
	TinyML models deployed on the ESP32 perform on-device inference for five tasks:
	
	\begin{enumerate}
		\item \textbf{Fall Detection} --- Identifies sudden motion changes
		(\texttt{FALL}/\texttt{NORMAL}) using MPU6050 accelerometer and
		gyroscope data.
		\item \textbf{Air Quality Classification} --- Categorizes conditions as
		\texttt{SAFE}, \texttt{SMOKE}, \texttt{TOXIC}, or \texttt{CO\_DANGER}
		using MQ sensor fusion.
		\item \textbf{Fire Validation} --- Combines flame sensor, temperature, and
		gas sensor outputs through multi-sensor fusion to reduce false positives.
		\item \textbf{Activity Recognition} --- Classifies
		\texttt{STILL}, \texttt{WALKING}, or \texttt{RUNNING} states.
		\item \textbf{Anomaly Detection} --- Detects unusual sensor patterns.
	\end{enumerate}
	
	This edge-based inference approach reduces latency and bandwidth usage by
	transmitting only processed results rather than raw sensor streams.
	
	% ─────────────────────────────────────────────────────────────────────────────
	%  V. DATA FORMAT
	% ─────────────────────────────────────────────────────────────────────────────
	\section{Data Format and Transmission}
	
	Each firefighter node transmits structured JSON data containing all monitored
	parameters:
	
	\begin{lstlisting}
		{
			"id": 1,
			"temp": 32,
			"mq2": 300,
			"mq9": 200,
			"mq135": 400,
			"lat": 17.12,
			"lon": 78.12,
			"status": "DANGER",
			"fall_status": "FALL",
			"air_quality": "TOXIC",
			"fire_valid": true,
			"activity": "RUNNING",
			"anomaly": false,
			"emergency": true,
			"timestamp": "ISO STRING"
		}
	\end{lstlisting}
	
	The backend safely parses incoming JSON packets from the LoRa serial stream,
	updates in-memory node tracking (\texttt{nodes = \{id: latest\_data\}}), emits
	Socket.io events for real-time frontend updates, and persists data to JSON
	session logs, CSV files, and an SQLite database for historical analysis and
	anomaly detection.
	
	% ─────────────────────────────────────────────────────────────────────────────
	%  VI. EXPERIMENTAL SETUP & RESULTS
	% ─────────────────────────────────────────────────────────────────────────────
	\section{Experimental Setup and Results}
	
	\subsection{Setup}
	
	\begin{itemize}
		\item \textbf{Sampling:} 10--20 iterations per condition
		\item \textbf{Technique:} Real-time continuous data collection
		\item \textbf{Environment:} Indoor and outdoor testing
		\item \textbf{Tools:} Arduino IDE, Serial Monitor, LoRa Library,
		TinyGPS++, Breadboard + Power Supply
	\end{itemize}
	
	\subsection{Procedure}
	
	\begin{enumerate}
		\item Connect all sensors to ESP32
		\item Initialize GPS, LoRa, and sensors
		\item Collect real-time data (temperature, gas, heart rate, location)
		\item Transmit data using LoRa
		\item Trigger alerts when thresholds are exceeded
	\end{enumerate}
	
	\subsection{Data Collection}
	
	\begin{table}[h]
		\caption{Observed Sensor Values}
		\label{tab:datacollection}
		\centering
		\begin{tabular}{|l|c|c|}
			\hline
			\textbf{Parameter} & \textbf{Observed Value} & \textbf{Status} \\
			\hline
			Temperature   & 45$^\circ$C           & High \\
			\hline
			Gas Level     & 300 ppm               & Moderate \\
			\hline
			Heart Rate    & 110 bpm               & Elevated \\
			\hline
			GPS           & Valid coordinates     & Active \\
			\hline
		\end{tabular}
	\end{table}
	
	\subsection{Trial Results}
	
	\begin{table}[h]
		\caption{Trial Results Summary}
		\label{tab:trials}
		\centering
		\begin{tabular}{|c|c|c|c|c|}
			\hline
			\textbf{Trial} & \textbf{Temp ($^\circ$C)} & \textbf{Gas (ppm)} &
			\textbf{HR (bpm)} & \textbf{Alert} \\
			\hline
			1 & 30 & 100 & 80  & No \\
			\hline
			2 & 45 & 300 & 110 & Yes \\
			\hline
			3 & 50 & 400 & 120 & Yes \\
			\hline
		\end{tabular}
	\end{table}
	
	The results demonstrate that average temperature increased progressively in
	simulated hazardous conditions. Alerts were consistently triggered when sensor
	thresholds were exceeded. Cross-validation of sensor readings against expected
	ranges and GPS output verification against known locations confirmed system
	accuracy. LoRa transmission was verified across multiple range tests, confirming
	reliable communication in simulated disaster conditions.
	
	% ─────────────────────────────────────────────────────────────────────────────
	%  VII. COMPARISON
	% ─────────────────────────────────────────────────────────────────────────────
	\section{Comparison with Existing Systems}
	
	\begin{table}[h]
		\caption{Proposed System vs.\ Existing Systems}
		\label{tab:comparison}
		\centering
		\renewcommand{\arraystretch}{1.2}
		\begin{tabular}{|p{1.8cm}|p{2.6cm}|p{2.6cm}|}
			\hline
			\textbf{Feature} & \textbf{Existing Systems} & \textbf{Proposed System} \\
			\hline
			Communication  & Short-range (Wi-Fi / Zigbee) & Long-range LoRa ($>$1~km) \\
			\hline
			Monitoring     & Partial (health OR location) & Multi-parameter (health + location + gas) \\
			\hline
			Alerts         & Limited or absent            & Real-time threshold-triggered \\
			\hline
			AI Inference   & Cloud-based or absent        & On-device TinyML (edge) \\
			\hline
			Data Storage   & Single format                & JSON, CSV, SQLite \\
			\hline
			Scalability    & Limited nodes                & Multi-node architecture \\
			\hline
		\end{tabular}
	\end{table}
	
	% ─────────────────────────────────────────────────────────────────────────────
	%  VIII. REAL-TIME DASHBOARD
	% ─────────────────────────────────────────────────────────────────────────────
	\section{Real-Time Dashboard}
	
	The frontend dashboard (TELEMETRY\_SYSTEM v2.4) is built using web technologies
	and provides multi-node visualization in real time. Key features include:
	
	\begin{itemize}
		\item \textbf{Territory Tracker:} Map-based visualization using Leaflet.js
		with dynamic markers per firefighter node. Color-coded status:
		Green (Safe), Red (Danger), Flashing Red (Emergency).
		\item \textbf{Active Entities Panel:} Dynamic cards per node showing
		temperature, gas levels, activity status, fall detection, and air quality.
		\item \textbf{Telemetry Streams:} Optional Chart.js graphs for temperature,
		gas concentration trends, and historical analysis.
		\item \textbf{Emergency Visualization:} Flashing red indicators, pulse
		animation, alert sound, and auto-focus on the affected node.
	\end{itemize}
	
	The UI follows a dark theme with glassmorphism styling, neon accents (cyan/red),
	and smooth animations. When a physical emergency button on a node is pressed,
	the emergency flag is set to \texttt{true}, GPS location is immediately
	transmitted via multiple redundant packets, and the frontend triggers full visual
	and audio alerts with auto-focus on the map.
	
	\subsection{Interactive Node Dashboard}
	
	Each firefighter node is represented as a dynamic card UI displaying temperature,
	gas levels, activity status, fall detection, and air quality. This aligns with
	modern IoT dashboards that emphasize clarity, modularity, and usability.
	
	\subsection{Backend Design}
	
	The backend is implemented using Node.js, Express, and Socket.io, with:
	\begin{itemize}
		\item Serial communication via \texttt{serialport}
		\item Safe JSON parsing for malformed-packet robustness
		\item In-memory node tracking: \texttt{nodes = \{id: latest\_data\}}
		\item Real-time WebSocket updates via Socket.io
		\item Multi-format storage: JSON logs, CSV files, SQLite database
	\end{itemize}
	
	% ─────────────────────────────────────────────────────────────────────────────
	%  IX. SDG ALIGNMENT
	% ─────────────────────────────────────────────────────────────────────────────
	\section{SDG Alignment}
	
	\textbf{SDG 3 -- Good Health and Well-Being:} The system continuously monitors
	vital parameters such as heart rate, body temperature, and gas exposure,
	providing real-time alerts that enable quicker response and timely rescue,
	significantly improving firefighter well-being and survival chances.
	
	\textbf{SDG 9 -- Industry, Innovation and Infrastructure:} The project
	incorporates innovative technologies such as IoT, ESP32 microcontrollers, and
	LoRa communication to develop reliable and resilient infrastructure for emergency
	response, especially where traditional communication fails.
	
	\textbf{SDG 11 -- Sustainable Cities and Communities:} The system enhances
	urban safety and disaster management by improving the efficiency and safety of
	firefighting operations, protecting lives, property, and infrastructure to make
	cities more resilient and sustainable.
	
	\textbf{SDG 13 -- Climate Action:} Firefighters play a crucial role in managing
	climate-related disasters such as wildfires. The system enhances their
	effectiveness and safety during such events, contributing to better climate
	disaster response and mitigation.
	
	% ─────────────────────────────────────────────────────────────────────────────
	%  X. ADVANTAGES
	% ─────────────────────────────────────────────────────────────────────────────
	\section{Advantages}
	
	\begin{itemize}
		\item Low latency due to on-device edge processing with TinyML.
		\item Scalable multi-node architecture supporting hundreds of firefighter nodes.
		\item Efficient long-range LoRa communication ($>$1~km) without infrastructure
		dependency.
		\item Robust emergency alert system with physical button, visual, and audio
		notifications.
		\item Multi-format data storage (JSON, CSV, SQLite) for analysis and
		reliability.
		\item Modular design enabling easy integration of additional sensors or cloud
		services.
	\end{itemize}
	
	% ─────────────────────────────────────────────────────────────────────────────
	%  XI. CONCLUSION
	% ─────────────────────────────────────────────────────────────────────────────
	\section{Conclusion}
	
	This paper presents a comprehensive IoT-based firefighter monitoring system
	integrating edge AI, LoRa communication, and real-time visualization. The system
	successfully demonstrated reliable multi-node communication via LoRa, accurate
	edge-based inference using TinyML, real-time monitoring with minimal latency,
	and effective emergency detection and visualization. The modular architecture
	allows easy scalability and integration with additional sensors or cloud
	services.
	
	Future work includes cloud integration for wider deployment, predictive analytics
	using historical data, enhanced AI models for improved decision-making, and
	exploration of indoor positioning alternatives to GPS for fully enclosed fire
	scenarios.
	
	% ─────────────────────────────────────────────────────────────────────────────
	%  REFERENCES
	% ─────────────────────────────────────────────────────────────────────────────
	\begin{thebibliography}{17}
		
		\bibitem{chaplia2023}
		O. Chaplia and H. Klym, ``Node.js framework for automated code scaling and
		execution on multiple distributed machines,'' \textit{2023 IEEE 18th
			International Conference on Computer Science and Information Technologies
			(CSIT)}, 2023.
		
		\bibitem{huang2020}
		X. Huang, ``Research and application of Node.js core technology,''
		\textit{2020 International Conference on Intelligent Computing and
			Human-Computer Interaction (ICHCI)}, 2020.
		
		\bibitem{senivongse2025}
		Senivongse and Kunphai, ``Model-driven development of web application backend
		using Node.js framework,'' 2025.
		
		\bibitem{haif2024}
		Haif \textit{et al.}, ``A decentralized dynamic relaying-based framework for
		enhancing LoRa networks performance,'' 2024.
		
		\bibitem{obrien2025}
		O'Brien \textit{et al.}, ``Design of a wireless cyber-physical system for gas
		leak detection with LoRa,'' 2025.
		
		\bibitem{abdelli2025}
		Abdelli \textit{et al.}, ``Field validated hybrid ESP-NOW and long range IoT
		monitoring system for energy autonomous precision agriculture,'' 2025.
		
		\bibitem{sole2022}
		Sol\'{e} \textit{et al.}, ``Implementation of a LoRa mesh library,'' 2022.
		
		\bibitem{deassis2025}
		de Assis \textit{et al.}, ``Scalable data acquisition system for real-time
		monitoring of photovoltaic plants,'' 2025.
		
		\bibitem{kumar2022}
		``Real-time GPS location tracking system using ESP32 and Neo-6M for IoT
		applications,'' 2023.
		
		\bibitem{mukherjee2022}
		``A novel machine learning based fall detection system using accelerometer and
		gyroscope,'' 2022.
		
		\bibitem{doddabasappla2021}
		``Statistical and machine learning-based recognition of coughing events using
		triaxial accelerometer sensor data,'' 2021.
		
		\bibitem{wang2023}
		``Evaluation of suitability of low-cost gas sensors for monitoring indoor and
		outdoor urban areas,'' 2020.
		
		\bibitem{wang2024}
		``Wearable low-cost and low-energy consumption gas sensor with machine learning
		to recognize outdoor areas,'' 2021.
		
		\bibitem{bolimera2025}
		``IoT based fire and motion detection system,'' 2022.
		
		\bibitem{ahmad2021}
		``On the evaluation of DHT22 temperature sensor for IoT application,'' 2021.
		
		\bibitem{yadav2026}
		``An IoT-enabled smart sensor system for real-time monitoring of soil and air
		conditions with integrated web-based farming guidance,'' 2022.
		
		\bibitem{dey2023}
		``Design and development of a smart and multipurpose IoT embedded system device
		using ESP32 microcontroller,'' 2023.
		
	\end{thebibliography}
	
\end{document}